\chapter{منابع و مراجع پیشنهادی برای مطالعه}
\label{app:references}

این طرح بر پایه‌ی تلفیق نظریات کلاسیک دموکراتیزاسیون و تجربه‌های عملی گذارهای دهه‌های اخیر تدوین شده است.

\section{متون کلیدی عدالت انتقالی و گذار}

\begin{itemize}
    \item \textit{Transitional Justice: How Emerging Democracies Reckon with Former Regimes}, Neil Kritz (ed.) — USIP, 1995.
    \item \textit{The Third Wave: Democratization in the Late Twentieth Century}, Samuel Huntington, 1991.
    \item \textit{Transitions from Authoritarian Rule}, O'Donnell, Schmitter \& Whitehead, 1986.
    \item \textit{The Spirit of Democracy}, Larry Diamond, 2008.
    \item \textit{No Future Without Forgiveness}, Desmond Tutu, 1999 (Experience of South Africa's TRC).
\end{itemize}

\section{منابع و نهادهای راهنما}

\begin{itemize}
    \item \textbf{International IDEA:} مرجع طراحی سیستم‌های انتخاباتی و دموکراسی (idea.int).
    \item \textbf{ICTJ:} مرکز بین‌المللی برای عدالت انتقالی (ictj.org).
    \item \textbf{Venice Commission:} نهاد مشورتی شورای اروپا در امور حقوق اساسی.
    \item \textbf{Constitute Project:} بانک جامع قوانین اساسی جهان (constituteproject.org).
    \item \textbf{V-Dem Institute:} شاخص‌های دموکراسی و حکمرانی (v-dem.net).
\end{itemize}

\section{منابع تکمیلی به زبان فارسی}

\begin{itemize}
    \item \textit{گذار به دموکراسی}، گزیده مقالات نظری و تطبیقی، ترجمه گروه‌های مختلف.
    \item \textit{عدالت در دوران گذار}، گزارش‌های مرکز اسناد حقوق بشر ایران.
    \item آرشیو بیانیه‌ها و منشورهای حقوقی گروه‌های مختلف اپوزیسیون ایران.
\end{itemize}

\vfill
\begin{center}
\small
\textit{«آینده‌ی ایران را نه با کینه از گذشته، بلکه با خرد برای آیندگان بسازیم.»}
\end{center}
